
\clearpage\part{Unbounded-density identities and the remaining radial interaction}


\section{Capacity reduction and exact weighted cancellations}
\label{sec:web-unbounded-reductions}
\hypertarget{proof:web-unbounded-reductions:s1}{%
\subsection{Inputs that do not use opposite-Weyl energy}\label{proof:web-unbounded-reductions:s1}}

All integrals below are normalized by the actual volume. The inherited strong-current and direct-resolvent arguments give, along a putative unrestricted zero-end sequence,

\[
 |v|\le C,
 \qquad |Lv|\le C,
 \qquad Lv-2v\longrightarrow0\quad\hbox{in }L^2,
 \qquad m=|v|^2\longrightarrow m_*={128\over81}\quad\hbox{in }L^2,
 \tag{\thesection.1.1}\label{proof:web-unbounded-reductions:e1.1}
\]

where \(L=\nabla^*\nabla\) on \(S^-\otimes L\). In particular

\[
 \fint|\nabla v|^2=\fint\operatorname{Re}\langle v,Lv\rangle
 \longrightarrow2m_*.
 \tag{\thesection.1.2}\label{proof:web-unbounded-reductions:e1.2}
\]

The positive spinor has \(\rho=|\Phi|^2\to\rho_*=4/3\), with \(d\rho\to0\) in \(L^2\). Its first torsion projection transfers to \(N_v\to0\) in \(L^2\) on each fixed positive-\(\rho\) region. None of these assertions requires a bound for \(\fint|W^-|^2\).

By contrast, the available bounds for \(\nabla^2v\) in \(L^2\) and \(\nabla v\) in \(L^4\) still depend on opposite-Weyl energy. The arguments below do not use those bounds.

\hypertarget{proof:web-unbounded-reductions:s2}{%
\subsection{Original-spinor zero regions can be removed by capacity, without an L4 estimate}\label{proof:web-unbounded-reductions:s2}}

The following elementary Caccioppoli estimate is useful here precisely because it only requires \eqref{proof:web-unbounded-reductions:e1.1}. For every smooth scalar cutoff \(0\le\chi\le1\), testing the equation for \(v\) by \(\chi^2v\) gives

\[
 \fint\chi^2|\nabla v|^2
 =\fint\chi^2\operatorname{Re}\langle Lv,v\rangle
 -2\fint\chi\operatorname{Re}\langle\nabla v,d\chi\otimes v\rangle.
\]

Young's inequality and the pointwise bounds in \eqref{proof:web-unbounded-reductions:e1.1} imply

\[
 \boxed{\fint\chi^2|\nabla v|^2
 \le C\fint\bigl(\chi^2+|d\chi|^2\bigr).}
 \tag{\thesection.2.1}\label{proof:web-unbounded-reductions:e2.1}
\]

Fix \(0<a<b<\rho_*\), and take \(\chi=\chi(\rho)\), equal to one on \(\rho\le a\) and zero on \(\rho\ge b\). Since \(\rho\to\rho_*\) in measure and \(d\rho\to0\) in \(L^2\),

\[
 \fint\chi^2+|d\chi|^2\longrightarrow0.
\]

Consequently

\[
 \boxed{\fint_{\{\rho\le a\}}|\nabla v|^2\longrightarrow0.}
 \tag{\thesection.2.2}\label{proof:web-unbounded-reductions:e2.2}
\]

Thus removing the original positive-spinor low-amplitude region does not need the finite-density L4 estimate. The remaining difficulty lies in low amplitude of \(v\) itself, where no corresponding capacity estimate has yet been obtained.

More generally, \eqref{proof:web-unbounded-reductions:e2.1} removes any exceptional set admitting scalar cutoffs whose normalized \(H^1\) norms tend to zero. Vanishing relative volume alone is insufficient for this purpose.

\hypertarget{proof:web-unbounded-reductions:s3}{%
\subsection{Any unresolved amplitude-gradient concentration is confined to lower wells}\label{proof:web-unbounded-reductions:s3}}

Use the positive scalar Laplacian \(\Delta=d^*d\). The exact identity is

\[
 \Delta m=2\operatorname{Re}\langle Lv,v\rangle-2|\nabla v|^2\le C.
 \tag{\thesection.3.1}\label{proof:web-unbounded-reductions:e3.1}
\]

For \(a\ge0\), use the Lipschitz test \((m-a)_+\), justified by smooth convex approximation. It gives

\[
 \fint_{\{m>a\}}|dm|^2
 =\fint(m-a)_+\Delta m
 \le C\fint(m-a)_+.
 \tag{\thesection.3.2}\label{proof:web-unbounded-reductions:e3.2}
\]

In particular,

\[
 \fint_{\{m>m_*\}}|dm|^2\longrightarrow0,
 \tag{\thesection.3.3}\label{proof:web-unbounded-reductions:e3.3}
\]

and for each \(0<\delta<m_*\),

\[
 \boxed{\limsup_j\fint_{\{m>m_*-\delta\}}|dm|^2\le C\delta.}
 \tag{\thesection.3.4}\label{proof:web-unbounded-reductions:e3.4}
\]

The first-order projection has the algebraic phase identity, with a sign depending on the chosen projection convention,

\[
 |j_v\mp J_+dm/2|\le C|v|\,|N_v|,
 \qquad j_v=\operatorname{Im}\langle v,\nabla v\rangle.
 \tag{\thesection.3.5}\label{proof:web-unbounded-reductions:e3.5}
\]

Equations \eqref{proof:web-unbounded-reductions:e2.2}, \eqref{proof:web-unbounded-reductions:e3.4}, and \eqref{proof:web-unbounded-reductions:e3.5} imply the analogous double-limit vanishing of \(j_v\) outside the lower-amplitude wells. Therefore the unresolved phase/radial defect cannot be an upper-amplitude spike. It is confined to sets on which \(m\) falls below its limiting value, whose relative volumes vanish for every fixed depth.

This conclusion is weaker than \(dm\to0\) in \(L^2\). In particular, \eqref{proof:web-unbounded-reductions:e3.4} does not bound the gradient energy on \(\{m\le m_*-\delta\}\). It would be incorrect to remove that set using its vanishing volume and the available L2 gradient bound.

\hypertarget{proof:web-unbounded-reductions:s4}{%
\subsection{The proposed weighted projection shortcut cancels exactly}\label{proof:web-unbounded-reductions:s4}}

A natural possible shortcut is to test the nearly vanishing holomorphic-type projection by \(v f(m)\), hoping to obtain a new positive \(f'(m)|dm|^2\) term without differentiating \(v\) twice. The relevant term cancels against the phase current.

Here is the pointwise algebra with fixed conventions. Take the Hermitian metric linear in its second argument and the projection

\[
 \mathsf P={I-iJ_+\over2},\qquad N=\mathsf P\nabla v.
\]

Then

\[
 \langle v,\nabla v\rangle={dm\over2}+ij_v,
 \qquad N\simeq0\quad\Longrightarrow\quad j_v\simeq J_+dm/2.
 \tag{\thesection.4.1}\label{proof:web-unbounded-reductions:e4.1}
\]

For real \(f=f(m)\),

\[
 N(fv)=fN+\mathsf P(df)\otimes v.
\]

Since \(\mathsf P\) is an orthogonal complex projection,

\[
 \begin{aligned}
 \operatorname{Re}\langle N,\mathsf P(df)\otimes v\rangle
 &=\frac14\langle dm,df\rangle
   -\frac12\langle j_v,J_+df\rangle\\
 &=\frac14f'(m)|dm|^2
   -\frac12\langle J_+dm/2,J_+f'(m)dm\rangle
   +\text{projection error}\\
 &=\text{projection error}.
 \end{aligned}
 \tag{\thesection.4.2}\label{proof:web-unbounded-reductions:e4.2}
\]

For fixed bounded \(f'\), the integral of that error tends to zero by \(N\to0\), bounded \(v\), and the L2 bound for \(dm\). Reversing the projection convention reverses the phase identity and the skew term together, leaving the cancellation unchanged.

The associated projection Weitzenboeck formula only involves the self-dual curvature contraction, which is attractive because that contraction contains no uncontrolled \(W^-\). However, its rough-Laplacian and curvature-potential terms already match at this endpoint. After the exact cancellation \eqref{proof:web-unbounded-reductions:e4.2}, it supplies no independent positive radial-gradient term. A proof using this test must provide an additional term not present in this formal weighted identity.

\hypertarget{proof:web-unbounded-reductions:s5}{%
\subsection{What the new L5 argument still requires in the unrestricted case}\label{proof:web-unbounded-reductions:s5}}

The positive line-bundle resolvent test at the end of the supplied fifth-line proof depends only on Ricci and determinant-line curvature bounds. That final inverse is not where the uncontrolled opposite-Weyl energy enters.

The unproved unrestricted connections occur earlier:

\begin{enumerate}
\def\labelenumi{\arabic{enumi}.}
\tightlist
\item
  The finite-density argument turns \(m\to m_*\) into \(dm\to0\). The above calculations remove the original-spinor zero region and locate the remaining radial defect, but do not eliminate the lower wells of \(m\).
\item
  The reduction to a single surviving off-diagonal coefficient requires the phase/radial error to be negligible. It is presently established globally by the finite-density gradient/Hessian estimates.
\item
  Even after that reduction, the weak CR errors contain products of a small L2 coefficient with \(\nabla v\). Without additional integrability this is only L1-small, whereas the line-resolvent test has an L4 first derivative, not a uniform pointwise first-derivative bound.
\item
  A positive L2 mass of the resulting line section need not survive a fixed amplitude truncation without uniform integrability. The finite-density L4 estimate supplies exactly that missing property.
\end{enumerate}

These issues cannot be removed merely by noting that the final test line bundle has no \(W^-\) curvature. Nor does \(\forall B<\infty\,\exists c(B)>0\) yield one positive constant uniform in \(B\).

\hypertarget{proof:web-unbounded-reductions:s6}{%
\subsection{Full Cartan calculation: a second attempted radial estimate also cancels}\label{proof:web-unbounded-reductions:s6}}

It is essential not to write \(b=c\theta_2\) before the radial current has been controlled. With \(N_v=Dv=0\) but with \(dm\) retained, the adapted-frame equations give

\[
 q={1\over2m}(m_2,-m_1,m_4,-m_3),
 \qquad
 b=d\theta_1+c\theta_2,
 \qquad
 d=-{m_3+im_4\over2m}.
 \tag{\thesection.6.1}\label{proof:web-unbounded-reductions:e6.1}
\]

Here \(m_i=e_i(m)\), and \(d\) in \eqref{proof:web-unbounded-reductions:e6.1} denotes a coefficient, not the exterior derivative. The additional coefficient follows from the second component of the Dirac equation: the radial component of \(\nabla_{\bar Z_2}v\) is cancelled by the second component of \(\nabla_{\bar Z_1}v\).

The complete Cartan equation is

\[
 d\theta_1=-i(\mathfrak a+p-q)\wedge\theta_1
             +\kappa\Xi_+-d\overline{\Xi_+},
 \qquad
 (F^{S^+})_{21}
 =(d\kappa+i\kappa(q-3p))\wedge\theta_1
       +\kappa^2\Xi_+-\kappa d\overline{\Xi_+}.
 \tag{\thesection.6.2}\label{proof:web-unbounded-reductions:e6.2}
\]

The coefficient \(c\) cancels from \(d\theta_1\), but \(d\) does not. Since

\[
 \kappa={\sqrt m\over2\rho^{3/2}},
 \qquad q\simeq-{J_+dm\over2m},
\]

the mass-dependent part in \eqref{proof:web-unbounded-reductions:e6.2} is

\[
 {\kappa\over2m}(dm-iJ_+dm)\wedge\theta_1.
 \tag{\thesection.6.3}\label{proof:web-unbounded-reductions:e6.3}
\]

Using \(\overline{\Xi_+}=\theta_1\wedge\theta_2\), \eqref{proof:web-unbounded-reductions:e6.3} equals

\[
 -\kappa{m_3+im_4\over2m}\,\overline{\Xi_+}
 =\kappa d\overline{\Xi_+}.
\]

It therefore cancels the last term in \eqref{proof:web-unbounded-reductions:e6.2} exactly. The first-plane component was already in the wedge kernel. Thus the full positive-connection equation does not supply the hoped-for second-plane radial estimate. Dropping \(d\theta_1\) from \(b\) at this stage would introduce a circular argument.

This does not invalidate the finite-density proof's reduction to \(b=c\theta_2+o_{L^2}(1)\): in that proof \(dm\to0\) has already been established, so the additional coefficient \(d\) is small. The issue concerns only an attempted extension that omits that established input.

There is nevertheless a potentially useful exact weight. If the second-plane component of \(dm\) can be made negligible at the needed strength, then \(J_+dm\simeq J_-dm\). The symbol from the fifth-line proof obeys

\[
 \mathcal L(iJ_-\alpha)=-\mathcal L(\alpha).
\]

Consequently, in the exact model with \(q=-J_-dm/(2m)\),

\[
 \boxed{
 \mathcal L(D_5c-3iqc)
 =m^{3/2}\mathcal L\bigl(D_5(m^{-3/2}c)\bigr).
 }
 \tag{\thesection.6.4}\label{proof:web-unbounded-reductions:e6.4}
\]

This conditional algebraic weight does not resolve the full system: the second-plane reduction is not provided by \eqref{proof:web-unbounded-reductions:e6.2}, and the full equation for \(b\) has the extra \(d\theta_1\) component and its derivatives. Formula \eqref{proof:web-unbounded-reductions:e6.4} must not be substituted for that full equation. Control at low \(m\) and a legal weak test also remain necessary.


\section{Double weighted tests and bilinear Hessian cancellation}
\label{sec:web-double-weighted}
\hypertarget{proof:web-double-weighted:s1}{%
\subsection{A weighted section estimate with no curvature term}\label{proof:web-double-weighted:s1}}

Let \(v\) be a smooth section of a Hermitian bundle with unitary connection and rough Laplacian \(L_v\). Suppose \[
\|v\|_\infty+\|L_vv\|_\infty\le C_v.
\] For every real scalar \(e\in H^1\), \[
\boxed{
\|e\nabla v\|_2
\le C(C_v)\bigl(\|e\|_2+\|de\|_2\bigr).
}
\tag{\thesection.1}\label{proof:web-double-weighted:e1}
\] To prove this, first take smooth \(e\), test the section equation against \(e^2v\), and integrate by parts: \[
\int e^2|\nabla v|^2
=\operatorname{Re}\int e^2\langle L_vv,v\rangle
-2\operatorname{Re}\int e\langle de\otimes v,\nabla v\rangle.
\] Young's inequality absorbs half of the left-hand side. No commutator, bundle curvature, Riemann tensor or Sobolev inequality occurs. Density proves the \(H^1\) version whenever the product is initially interpreted by approximation.

For the norm of a smooth bundle-valued tensor \(E\), use \(e_\delta=(\delta^2+|E|^2)^{1/2}\), Kato's inequality, and let \(\delta\downarrow0\). Hence \[
\boxed{
\||E|\nabla v\|_2\le C(C_v)\|E\|_{H^1}.
}
\tag{\thesection.2}\label{proof:web-double-weighted:e2}
\] The estimate requires no uniform bound for \(E\) itself.

\hypertarget{proof:web-double-weighted:s2}{%
\subsection{Two applications that replace the growing fourth-derivative budget}\label{proof:web-double-weighted:s2}}

Let \(w\) be a section of another Hermitian bundle with \[
\|w\|_\infty+\|L_ww\|_\infty\le C_w,
\qquad \|\nabla w\|_2+\|\nabla^2w\|_2\le C_{w,2}.
\] Let \(\Phi\) be a further section satisfying \[
\|\nabla\Phi\|_2+\|\nabla^2\Phi\|_2\le C_{\Phi,2}.
\] Applying \eqref{proof:web-double-weighted:e2} first to \(v\), with \(E=\nabla w\), and then to \(w\), with \(E=\nabla\Phi\), gives \[
\boxed{
\||\nabla v||\nabla w|\|_2
\le C(C_v)C_{w,2},
}
\tag{\thesection.3}\label{proof:web-double-weighted:e3}
\] \[
\boxed{
\||\nabla\Phi||\nabla w|\|_2
\le C(C_w)C_{\Phi,2}.
}
\tag{\thesection.4}\label{proof:web-double-weighted:e4}
\] These are not consequences of an \(L^4\) bound for \(\nabla v\). In particular their constants do not depend on \(\fint|W^-|^2\).

For the actual positive line resolvent, \[
(L_w+10/3)w=T_M(c),
\] the fixed height \(M\) yields bounded \(w,L_ww\). Its line curvature is \(5F_A\); the already established bounded-model Ricci and line-curvature estimates give \(C_{w,2}\) independent of the negative Weyl budget. Thus \eqref{proof:web-double-weighted:e3}–\eqref{proof:web-double-weighted:e4} apply to those tests once \(M\) is fixed.

\hypertarget{proof:web-double-weighted:s3}{%
\subsection{Bilinear Hessian cancellation by one integration by parts}\label{proof:web-double-weighted:s3}}

Let \(\zeta\) be a one-form or a section of a fixed-rank tensor bundle, with \[
\|\zeta\|_{H^1}\to0.
\] Let \(\mathcal C\) be a globally defined smooth contraction tensor of the appropriate bundle type. Assume uniformly \[
|\mathcal C|\le C_0,
\qquad
|\nabla\mathcal C|\le C_0(1+|\nabla\Phi|+|\nabla v|).
\tag{\thesection.5}\label{proof:web-double-weighted:e5}
\] For any fixed index pattern, including the symmetric negative Hessian, define \[
I=\operatorname{Re}\fint
\zeta_a\mathcal C^{iajk}\bigl(\nabla_i\nabla_jv,\nabla_kw\bigr).
\] Then \[
\boxed{|I|\le C\|\zeta\|_{H^1}\longrightarrow0.}
\tag{\thesection.6}\label{proof:web-double-weighted:e6}
\] The constant depends on the displayed bounds, but not on \(\|\nabla^2v\|_2\), \(\|\nabla v\|_4\), or the opposite Weyl curvature.

\textbf{Proof.} Integrate the \(i\) derivative off \(\nabla_jv\). Since all contractions are covariant and the manifold is closed, there are exactly three types of terms: \[
\begin{aligned}
I_1&=-\operatorname{Re}\fint(\nabla_i\zeta_a)\mathcal C^{iajk}
 (\nabla_jv,\nabla_kw),\\
I_2&=-\operatorname{Re}\fint\zeta_a(\nabla_i\mathcal C^{iajk})
 (\nabla_jv,\nabla_kw),\\
I_3&=-\operatorname{Re}\fint\zeta_a\mathcal C^{iajk}
 (\nabla_jv,\nabla_i\nabla_kw).
\end{aligned}
\] There is no interchange of covariant derivatives. Thus no curvature is introduced.

By \eqref{proof:web-double-weighted:e3}, \[
|I_1|\le C\|\nabla\zeta\|_2
\||\nabla v||\nabla w|\|_2
\le C\|\nabla\zeta\|_2.
\] Apply \eqref{proof:web-double-weighted:e2} with \(E=\zeta\) to obtain \[
\||\zeta|\nabla v\|_2\le C\|\zeta\|_{H^1}.
\tag{\thesection.7}\label{proof:web-double-weighted:e7}
\] Consequently \[
|I_3|\le C\||\zeta|\nabla v\|_2\|\nabla^2w\|_2
\le C\|\zeta\|_{H^1}.
\] For \(I_2\), use \eqref{proof:web-double-weighted:e5}, \eqref{proof:web-double-weighted:e3}, \eqref{proof:web-double-weighted:e4}, and \eqref{proof:web-double-weighted:e7}: \[
\begin{aligned}
|I_2|
&\le C\||\zeta|\nabla v\|_2
\left(\|\nabla w\|_2+
\||\nabla\Phi||\nabla w|\|_2+
\||\nabla v||\nabla w|\|_2\right)\\
&\le C\|\zeta\|_{H^1}.
\end{aligned}
\] This proves \eqref{proof:web-double-weighted:e6}. In particular, \textbf{no antisymmetry of \(i,j\) is necessary}. The same proof applies termwise to any fixed finite linear combination or to the symmetrized Hessian.

\hypertarget{proof:web-double-weighted:s4}{%
\subsection{Fixed amplitude cutoffs and adapted-frame tensors}\label{proof:web-double-weighted:s4}}

An expression written in adapted frames is eligible for \eqref{proof:web-double-weighted:e6} only after it is a genuine globally defined contraction. If it is built tensorially from bounded \(\Phi,v\), parallel Clifford multiplication, and smooth functions of their norms which vanish on fixed lower-amplitude neighborhoods, then its denominators are bounded wherever it is nonzero. Its covariant derivative satisfies \eqref{proof:web-double-weighted:e5}. Cutoff derivatives are included in \eqref{proof:web-double-weighted:e5}, since \[
|d\chi(|v|^2)|\le C|\nabla v|,
\qquad |d\chi(|\Phi|^2)|\le C|\nabla\Phi|.
\] This verifies eligibility for such tensors without assuming a global adapted frame. It does not automatically verify that any particular frame expression descends correctly; that algebraic check remains necessary.

\hypertarget{proof:web-double-weighted:s5}{%
\subsection{Further directly usable estimates}\label{proof:web-double-weighted:s5}}

The same argument immediately gives, for a tensor error \(R\), \[
\left|\fint R*(\nabla v)*(\nabla w)\right|
\le C\|R\|_2.
\tag{\thesection.8}\label{proof:web-double-weighted:e8}
\] Thus a small \(L^2\) first-order error can be paired with these two gradients without paying a separate \(\nabla v\) fourth-power norm. Also \[
\fint |\zeta||\nabla v|^2|\nabla w|
\le \||\zeta|\nabla v\|_2
     \||\nabla v||\nabla w|\|_2
\to0.
\tag{\thesection.9}\label{proof:web-double-weighted:e9}
\] These estimates identify a class of growing-budget losses in the earlier CR bookkeeping that can genuinely be removed.

\hypertarget{proof:web-double-weighted:s6}{%
\subsection{The remaining radial interaction}\label{proof:web-double-weighted:s6}}

The new lemma resolves the displayed bilinear term involving \(\zeta\nabla^2v\nabla w\), and also the small-error triple product \eqref{proof:web-double-weighted:e8}. It avoids the uncontrolled negative Hessian rather than merely assigning it a slower growth rate.

For a proof of unrestricted the weak endpoint problem, one must still derive the \textbf{complete} actual weak CR identity with every term accounted for and show it belongs to this class or another uniformly controlled class. In particular, this lemma does not prove \(d|v|^2\to0\), does not automatically give nonvanishing uniformly integrable mass for the fifth-line section, and does not make a radial coefficient small when its only information is a bounded \(H^1\) norm. Fixed amplitude cutoffs do not erase these requirements.

The fixed-height resolvent hypothesis also remains meaningful: if the mass of \(c\) escapes to arbitrarily high values, the existing contradiction requires a separate argument that fixes the truncation height. This lemma should therefore be inserted into the actual proof and the remaining radial and mass terms checked explicitly; it is not a license to declare the full endpoint solved.


\section{Amplitude defect measures and weak radial tests}
\label{sec:web-radial-measure}
\hypertarget{proof:web-radial-measure:s2}{%
\subsection{The exact defect law retains a positive radial measure}\label{proof:web-radial-measure:s2}}

The amplitudes lie in a fixed interval \([0,M]\). Define positive measures on that interval by \[
 \mu_j(B)=\fint_{\{m_j\in B\}}|dm_j|^2,
 \qquad \nu_j(B)=\fint_{\{m_j\in B\}}q_j.
\] Their total masses are uniformly bounded, by (1). Pass to weakly convergent subsequences \(\mu_j\rightharpoonup\mu\), \(\nu_j\rightharpoonup\nu\). For every \(f\in C^1([0,M])\), the scalar equation tested with \(f(m_j)\) gives \[
 \int f'(t)\,d\mu(t)
 =4m_*f(m_*)-2\int f(t)\,d\nu(t).
 \tag{\thesection.2}\label{proof:web-radial-measure:e2}
\] Consequently \[
 \mu'=2\nu-4m_*\delta_{m_*}
 \tag{\thesection.3}\label{proof:web-radial-measure:e3}
\] as distributions after extension by zero outside \([0,M]\). In particular \(\nu([0,M])=2m_*\). The upper-amplitude estimate from the preceding continuation shows that \(\mu\) is supported in \([0,m_*]\); it does not make \(\mu\) vanish there.

These statements are not merely a technical failure of one estimate. The measures \[
 d\mu(t)=4t\,\mathbf1_{(0,m_*)}(t)\,dt,
 \qquad d\nu(t)=2\,\mathbf1_{(0,m_*)}(t)\,dt
 \tag{\thesection.4}\label{proof:web-radial-measure:e4}
\] satisfy \eqref{proof:web-radial-measure:e2}–\eqref{proof:web-radial-measure:e3}, including the endpoint atom of the distributional derivative, and satisfy the Kato bound \(d\mu\le4t\,d\nu\). They even satisfy the stronger radial coefficient equality \(d\mu=2t\,d\nu\) occurring in the formal exact first-jet model with no surviving \(c\)-component. Thus the entire mean gradient energy can be carried by lower-amplitude radial concentration without contradicting these scalar identities.

Equation \eqref{proof:web-radial-measure:e4} is a measure-level consistency test, not an actual field construction. In particular it does not establish the full equations \(N_v\simeq0\), \(D^-v\simeq0\), or the Levi-Civita realization of the connection.

\hypertarget{proof:web-radial-measure:s3}{%
\subsection{A fixed flat scalar realization of the amplitude-concentration phenomenon}\label{proof:web-radial-measure:s3}}

Here is a more concrete check: all the scalar features of (1) except their realization as the covariant derivatives of an actual the weak endpoint problem spinor can occur on a fixed flat four-torus. This prevents the scalar-amplitude shortcut even under excellent base geometry.

Let \(\varphi\ge0\) be smooth, supported in a small Euclidean ball, with integral one. On the unit flat torus let \[
 \rho_\delta(y)=\delta^{-4}\varphi(y/\delta),\qquad
 \Delta U_\delta=\rho_\delta-1,\qquad\int U_\delta=0.
\] Write \(G\) for the mean-zero torus Green function. Locally \(G(y)=c_4|y|^{-2}+O(1)\), and away from its pole \(G\) is smooth. Hence, directly from convolution with \(\rho_\delta\), \[
 -C_0\le U_\delta\le C_1\delta^{-2},\qquad
 \int U_\delta^2\le C_2(1+|\log\delta|),
 \tag{\thesection.5}\label{proof:web-radial-measure:e5}
\] and \[
 \int|dU_\delta|^2=\int U_\delta\rho_\delta
 =c_\varphi\delta^{-2}+O(1),\qquad c_\varphi>0.
 \tag{\thesection.6}\label{proof:web-radial-measure:e6}
\] The last coefficient is the positive Euclidean Newtonian self-energy of \(\varphi\); the mean-zero condition removes the constant term in the source.

Take \(\varepsilon=1/n\), put \(\delta=\varepsilon\), and define an \(\varepsilon\)-periodic function on the unit torus by \[
 u_\varepsilon(x)=a\varepsilon^2\bigl(U_\varepsilon(x/\varepsilon)+C_0\bigr).
 \tag{\thesection.7}\label{proof:web-radial-measure:e7}
\] Choose the fixed \(a>0\) small enough that \(0\le u_\varepsilon\le1/4\) for all sufficiently small \(\varepsilon\), and also \(a\le1\). The scaling and \eqref{proof:web-radial-measure:e5}–\eqref{proof:web-radial-measure:e6} give \[
 u_\varepsilon\to0\text{ in }L^2,\quad
 \int|du_\varepsilon|^2\to a^2c_\varphi>0,\quad
 \Delta u_\varepsilon=a\bigl(\rho_\varepsilon(x/\varepsilon)-1\bigr)\ge-a.
 \tag{\thesection.8}\label{proof:web-radial-measure:e8}
\] Define \[
 m_\varepsilon=m_*(1-u_\varepsilon)^2,\qquad
 q_\varepsilon=2m_\varepsilon-\frac12\Delta m_\varepsilon.
 \tag{\thesection.9}\label{proof:web-radial-measure:e9}
\] For the positive Laplacian, \[
 \Delta m_\varepsilon
 =-2m_*(1-u_\varepsilon)\Delta u_\varepsilon
   -2m_*|du_\varepsilon|^2\le2am_*.
\] Moreover \[
 q_\varepsilon
 =2m_*(1-u_\varepsilon)^2
 +m_*(1-u_\varepsilon)\Delta u_\varepsilon
 +m_*|du_\varepsilon|^2
 \ge m_*|du_\varepsilon|^2
 =\frac{|dm_\varepsilon|^2}{4m_\varepsilon},
 \tag{\thesection.10}\label{proof:web-radial-measure:e10}
\] because \(2(1-u_\varepsilon)-a\ge1/2\). Thus \[
 m_\varepsilon\to m_*\text{ in }L^2,\quad
 \fint q_\varepsilon=2\fint m_\varepsilon\to2m_*,\quad
 \liminf\int|dm_\varepsilon|^2>0,
 \tag{\thesection.11}\label{proof:web-radial-measure:e11}
\] while \(m_\varepsilon\) is bounded above and uniformly bounded away from zero, and \(\Delta m_\varepsilon\le C\). The base is one fixed flat torus: no collapsing coordinates or unbounded Ricci curvature are used.

Each cell has diameter \(\varepsilon\), and the source core has physical radius \(\varepsilon^2\). There are \(\varepsilon^{-4}\) cells. The four-dimensional capacities of the cores therefore have order \[
 \varepsilon^{-4}(\varepsilon^2)^2\asymp1,
\] whereas their total volume has order \(\varepsilon^{-4}(\varepsilon^2)^4=\varepsilon^4\to0\). This accounts for the coexistence of nearly constant amplitude and nonvanishing radial energy.

\textbf{Scope:} \(q_\varepsilon\) is a nonnegative scalar energy density satisfying the amplitude equation and Kato constraint. This construction does not claim \(q_\varepsilon=|\nabla v_\varepsilon|^2\) for the required spinor and spin connection. It is not an the weak endpoint problem counterexample. Its purpose is to rule out deducing the missing radial estimate from scalar amplitude/Laplacian information alone.

\hypertarget{proof:web-radial-measure:s4}{%
\subsection{Information not supplied by the scalar model}\label{proof:web-radial-measure:s4}}

The inherited first projection and Dirac equation yield, on a positive-amplitude adapted frame and with their small errors retained, \[
 q_{\rm ph}=-\frac{J_+dm}{2m},\qquad
 b=d\theta_1+c\theta_2,\qquad
 d=-\frac{m_3+im_4}{2m}.
\] The exact positive-curvature Cartan equation cancels the displayed radial term; it does not imply \(d\to0\). This cancellation was verified in the preceding continuation and is not re-counted as a new result here.

A proposed nonintegrability route uses \([\nabla_{Z_1},\nabla_{Z_2}]v\), whose relevant curvature component contains Ricci rather than uncontrolled \(W^-\). The small first projection yields a weak differentiated constraint. Turning it into an \(L^2\) estimate for the missing derivative requires pairing a small first-order quantity with a derivative of the test. The currently available test involves \(\nabla^2v\) or products with \(\nabla v\); their control is precisely what grows with opposite-Weyl energy. Merely invoking the bracket identity or local exact almost-Kähler rigidity does not bound that pairing on the degenerating sequence.

No new estimate obtained in this attempt eliminates the measure \(\mu\), controls the required bracket test without the missing integrability, or constructs a closed actual escape. The precise remaining question is whether the full positive-spinor/Levi-Civita equations force this capacity defect to vanish or force it to pay a nonzero amount of the original the weak endpoint problem residual. A proof from the scalar consequences alone is excluded by subsection~\ref{proof:web-radial-measure:s3}.

\hypertarget{proof:web-radial-measure:s5}{%
\subsection{A genuine estimate from the full section equation: small coefficients multiply the gradient in L2}\label{proof:web-radial-measure:s5}}

There is additional uniform information, beyond the scalar defect law. Let \(e\) be any smooth section of any metric bundle (a scalar or a one-form is allowed), and let \(v\) be the actual bounded auxiliary section with \(|v|+|Lv|\le C\). Then \[
 \boxed{\fint |e|^2|\nabla v|^2
 \le C\fint\bigl(|e|^2+|\nabla e|^2\bigr).}
 \tag{\thesection.12}\label{proof:web-radial-measure:e12}
\] The constant depends only on the bounds for \(v,Lv\). It contains no opposite-Weyl bound and no Sobolev constant. Indeed test \(Lv\) by \(|e|^2v\): \[
 \int |e|^2|\nabla v|^2
 =\operatorname{Re}\int |e|^2\langle Lv,v\rangle
 -\operatorname{Re}\int \langle d|e|^2\otimes v,\nabla v\rangle.
\] Use \(|d|e|^2|\le2|e||\nabla e|\) and absorb half of the left side by Young. The argument is valid without an a priori pointwise bound on \(e\).

The inherited strong-current estimate has \(\zeta\to0\) in normalized \(H^1\). Thus the actual system gives \[
 \boxed{\|\zeta\,\nabla v\|_{2,\mathrm{av}}\to0.}
 \tag{\thesection.13}\label{proof:web-radial-measure:e13}
\] This is stronger than the naive product estimate \(L^4\cdot L^2\subset L^{4/3}\). Similarly, whenever an actual resolvent/cap error \(e\to0\) in \(H^1\), one obtains \(e\nabla v\to0\) in \(L^2\).

Consequently, for a line-bundle test section \(w\) with \(\|\nabla^2w\|_2\le C\), terms of the form \[
 \fint C_0*(\zeta\nabla v)*\nabla^2w
\] tend to zero whenever the tensor coefficient \(C_0\) is uniformly bounded. This removes one genuine weak-test product obstruction without controlling \(\nabla^2v\).

\hypertarget{proof:web-radial-measure:A-second-use-of-Caccioppoli-removes-the-previously-problematic-differentiated-current-tests}{%
\subsubsection{A second use of Caccioppoli removes the previously problematic differentiated-current tests}\label{proof:web-radial-measure:A-second-use-of-Caccioppoli-removes-the-previously-problematic-differentiated-current-tests}}

Suppose the line-bundle test satisfies \[
 |w|+|\nabla^*\nabla w|\le C,
 \qquad \|\nabla w\|_2+\|\nabla^2w\|_2\le C.
\] These are the existing bounded-forcing positive line-resolvent bounds, independent of opposite-Weyl energy. Apply \eqref{proof:web-radial-measure:e12} with \(e=\nabla w\), using Kato for the scalar norm. This gives the genuinely uniform product estimate \[
 \boxed{\|\,|\nabla v|\,|\nabla w|\,\|_2\le C.}
 \tag{\thesection.14}\label{proof:web-radial-measure:e14}
\] Although neither factor has a uniform pointwise bound, their product is controlled in \(L^2\). Similarly apply Caccioppoli to the bounded line section \(w\) with \(e=\nabla\Phi\). The inherited positive-spinor Hessian bound yields \[
 \boxed{\|\,|\nabla\Phi|\,|\nabla w|\,\|_2\le C.}
 \tag{\thesection.15}\label{proof:web-radial-measure:e15}
\]

For a bounded tensor coefficient \(C^{iajk}\) satisfying \[
 |\nabla C|\le C_0(1+|\nabla\Phi|+|\nabla v|),
\] integrating the representative differentiated-current pairing by parts gives \[
 \begin{aligned}
 \int C^{iajk}(\nabla_i\zeta_a)
                  \langle\nabla_jv,\nabla_kw\rangle
 ={}&-\int\zeta_a C^{iajk}
               \langle\nabla_i\nabla_jv,\nabla_kw\rangle\\
 &-\int\zeta_a C^{iajk}
               \langle\nabla_jv,\nabla_i\nabla_kw\rangle\\
 &-\int\zeta_a(\nabla_iC^{iajk})
               \langle\nabla_jv,\nabla_kw\rangle.
 \end{aligned}
 \tag{\thesection.16}\label{proof:web-radial-measure:e16}
\] The left side tends to zero by \(\|\nabla\zeta\|_2\to0\) and \eqref{proof:web-radial-measure:e14}. The second line tends to zero by \eqref{proof:web-radial-measure:e13} and the test Hessian bound. The third line tends to zero by \eqref{proof:web-radial-measure:e13} paired with \eqref{proof:web-radial-measure:e14}–\eqref{proof:web-radial-measure:e15}. Thus \textbf{the first line, including its symmetric Hessian component, tends to zero as well}. No alternating-index assertion or negative-spinor Hessian estimate is needed. All integrations are divided by the actual volume.

More explicitly, each such pairing is bounded by a fixed constant times \[
 \|\nabla\zeta\|_2+\|\zeta\|_{H^1},
\] with the displayed test/coefficient budgets fixed. If the current estimate is \(\|\zeta\|_{H^1}=O(x^\gamma)\), this argument keeps the same exponent; no factor involving \(b\) is introduced.

This repairs a real testing obstruction. It must not be overstated: the estimate for \(C\) is to be checked for the full tensor and every cutoff. A coefficient involving a fixed smooth amplitude cutoff has such a derivative bound on its support. It does not justify dividing by an uncontrolled amplitude without that cutoff.

\hypertarget{proof:web-radial-measure:Linear-weak-radial-tests-also-vanish--without-strong-radial-convergence}{%
\subsubsection{Linear weak radial tests also vanish, without strong radial convergence}\label{proof:web-radial-measure:Linear-weak-radial-tests-also-vanish--without-strong-radial-convergence}}

Let \(C=C(v,\Phi)\) be a uniformly bounded contraction tensor, with \[
 |\nabla C|\le C_0(1+|\nabla v|+|\nabla\Phi|).
\] Then, for every line-resolvent test \(w\) obeying the bounds above, \[
 \boxed{\left|\fint dm*C*\nabla w\right|
 \le C\|m-m_*\|_2\to0.}
 \tag{\thesection.17}\label{proof:web-radial-measure:e17}
\] Indeed integration by parts rewrites the integral as \[
 -\fint(m-m_*)\bigl(C*\nabla^2w+(\nabla C)*\nabla w\bigr).
\] The bracket is uniformly \(L^2\) by the test Hessian bound and \eqref{proof:web-radial-measure:e14}–\eqref{proof:web-radial-measure:e15}. The same conclusion permits smooth fixed lower-amplitude cutoffs in \(C\): derivatives of such cutoffs are bounded by fixed multiples of \(|\nabla v|+|\nabla\Phi|\), so their transition errors are already included in this calculation.

This is a useful weak radial cancellation, but the order of the test is essential. It does \textbf{not} justify discarding \[
 \fint dm*(\nabla v)*\nabla w
 \quad\hbox{or}\quad
 \fint dm*C*\nabla^2w.
 \tag{\thesection.18}\label{proof:web-radial-measure:e18}
\] Moving the derivative in the first expression introduces \((m-m_*)\nabla^2v\nabla w\); moving it in the second introduces third derivatives of \(w\). Neither has been controlled here. In particular \eqref{proof:web-radial-measure:e17} is not strong \(L^2\) convergence of \(dm\), and it does not yet remove the radial-square mass in the full Cartan system.

\hypertarget{proof:web-radial-measure:What-remains-after-this-repair}{%
\subsubsection{What remains after this repair}\label{proof:web-radial-measure:What-remains-after-this-repair}}

The unrestricted first-jet reduction still has the radial coefficient \[
 d=-\frac{m_3+i m_4}{2m}
\] and the phase current \(-J_+dm/(2m)\). Neither is known to be \(L^2\)-small. Formula \eqref{proof:web-radial-measure:e12} cannot be invoked with \(e=m-m_*\) as a small coefficient: that would require the unknown smallness of \(dm\). The surviving first-jet norm is \[
 |\nabla v|^2=\frac{E_1+2E_2}{2m}+2m|c|^2,
 \qquad E_1=m_1^2+m_2^2,\quad E_2=m_3^2+m_4^2.
\] Thus the fixed total gradient mass does not yet produce a fixed nonzero mass for the fifth-line coefficient \(c\). If this radial concentration is not first eliminated, the nonnegative contribution \((E_1+2E_2)/(2m)\) may consume the mass used by the finite-density contradiction.

Nor does the bounded-test estimate alone prevent \(|c|^2\) from concentrating above every fixed amplitude truncation. The general scalar constructions in subsection~\ref{proof:web-radial-measure:s3} are not solutions of the full field system; they show why these two remaining steps require that full system, instead of another scalar manipulation of (1).

The actual new input is therefore \eqref{proof:web-radial-measure:e13}–\eqref{proof:web-radial-measure:e16}, which removes differentiated-current/negative-frame test products from the list of obstructions under the stated tensor bounds. The amplitude-capacity defect and its interaction with the complete radial Cartan system remain unresolved.


\section{Strong radial transfer to the original actual field}
\label{sec:web-mass-transfer}
\hypertarget{proof:web-mass-transfer:s1}{%
\subsection{A strong radial transfer for bounded resolvent approximations}\label{proof:web-mass-transfer:s1}}

Suppose \(u,v\) are sections of the same Hermitian bundle on a closed manifold, with \[
\|u\|_\infty+\|v\|_\infty+\|\nabla^*\nabla v\|_\infty\le C,
\qquad h:=v-u.
\] Then, with normalized norms, \[
\boxed{
\big\||v|^2-|u|^2\big\|_{H^1}\le C'\|h\|_{H^1}.
}
\tag{\thesection.1}\label{proof:web-mass-transfer:e1}
\] The zero-order term is immediate. For the derivative write \[
d(|v|^2-|u|^2)
=2\operatorname{Re}\langle v,\nabla h\rangle
 +2\operatorname{Re}\langle h,\nabla v\rangle
 -2\operatorname{Re}\langle h,\nabla h\rangle.
\] The first and third terms are bounded by \(C\|\nabla h\|_2\). The middle term is controlled by the weighted Caccioppoli estimate from E: \[
\||h|\nabla v\|_2\le C\|h\|_{H^1}.
\] This proves \eqref{proof:web-mass-transfer:e1}, without a negative-spinor Hessian estimate or an opposite Weyl-energy budget.

Apply this to the bounded actual torsion cap \(u=b(|\tau|^2)\tau\) and its bounded-forcing resolvent \(v\). Since the existing construction gives \(u-v\to0\) in normalized \(H^1\), \[
d|v|^2-d|u|^2\to0\quad\text{in normalized }L^2.
\tag{\thesection.2}\label{proof:web-mass-transfer:e2}
\] Thus the unresolved radial derivative is already a derivative of a bounded smooth function of the \textbf{actual} torsion magnitude; it is not created by the resolvent approximation.

\hypertarget{proof:web-mass-transfer:s2}{%
\subsection{The original actual mass identifies the differential order exactly}\label{proof:web-mass-transfer:s2}}

Put \(\rho=|\Phi|^2\), \(q_\Phi=|\nabla_A\Phi|^2\), and \[
\nu=2\rho^2q_\Phi-|\tau|^2.
\] The exact scalar equation gives \[
\Delta\rho=-2q_\Phi+2\rho-\rho^2.
\] Therefore \[
\boxed{|\tau|^2=2\rho^3-\rho^4-\rho^2\Delta\rho-\nu.}
\tag{\thesection.3}\label{proof:web-mass-transfer:e3}
\] Differentiating, \[
d|\tau|^2
=(6\rho^2-4\rho^3-2\rho\Delta\rho)d\rho
 -\rho^2d\Delta\rho-d\nu.
\tag{\thesection.4}\label{proof:web-mass-transfer:e4}
\] The strong-current estimates control \(\Delta\rho\) in \(L^2\), equivalently the trace of \(\nabla\zeta\); they do not yet control \(d\Delta\rho\) in \(L^2\). Substituting \eqref{proof:web-mass-transfer:e4} into a radial term and integrating by parts moves \textbf{two derivatives of the current} onto the rest of the test, unless the complete combination has an additional cancellation. This issue must be handled in the combined actual equation, not by treating \(d|\tau|^2\) as an already small first-order error.

Equation \eqref{proof:web-mass-transfer:e2} means that such a successful combined cancellation may be proved using the actual identity \eqref{proof:web-mass-transfer:e3}, then transferred to the resolvent field. It avoids losing a factor of the negative-curvature budget merely in that transfer.

\hypertarget{proof:web-mass-transfer:s3}{%
\subsection{A one-sided radial localization that follows without density control}\label{proof:web-mass-transfer:s3}}

Let \(a=|v|^2\), with \(a\to m_*>0\) in normalized \(L^2\), \(v\) uniformly bounded, and \[
\nabla^*\nabla v=2v+r_v,\qquad \|r_v\|_2\to0.
\] Then \[
\Delta a=4a-2|\nabla v|^2+e_a,
\qquad\|e_a\|_2\to0.
\] Testing with \((a-m_*)_+\) yields \[
\int_{a>m_*}|da|^2
\le 4\int(a-m_*)_+a+\int(a-m_*)_+e_a\to0.
\tag{\thesection.5}\label{proof:web-mass-transfer:e5}
\] For \(0<\delta<m_*\), test instead with \((a-m_*+\delta)_+\). Dropping the nonpositive gradient-energy term gives \[
\limsup_j\fint_{a_j>m_*-\delta}|da_j|^2
\le 4m_*\delta.
\tag{\thesection.6}\label{proof:web-mass-transfer:e6}
\] Thus any persistent radial energy must be supported increasingly on genuinely lower-amplitude regions. It cannot be placed exclusively in small fluctuations above the limiting amplitude.

This does not remove the defect: those lower-amplitude regions may have vanishing volume and nonvanishing radial energy. It also does not supply uniform integrability of \(|\nabla v|^2\) or nonzero mass of the fifth-line horizontal component. The scalar amplitude equation permits precisely the concentration alternative that the unrestricted proof must decide.

\hypertarget{proof:web-mass-transfer:s4}{%
\subsection{The remaining combined equation}\label{proof:web-mass-transfer:s4}}

The double-weighted lemma in E controls \(\zeta\nabla^2v\nabla w\), and a small \(L^2\) first-order coefficient paired with \(\nabla v\nabla w\), uniformly in the opposite curvature. It does not alone control a term with \(d\Delta\rho\) paired with a tensor containing \(\nabla^2w\): direct integration by parts differentiates the test a third time. A proof that the full radial CR combination cancels that derivative has not been obtained here.

No new uncontrolled quantity has been hidden in the notation: \eqref{proof:web-mass-transfer:e3} gives the exact actual-field source of the radial defect, \eqref{proof:web-mass-transfer:e1} removes the resolvent-transfer loss, and \eqref{proof:web-mass-transfer:e5}–\eqref{proof:web-mass-transfer:e6} locate where any remaining concentration can occur. The final density-free implication remains open.


\section{The full polynomial radial one-form}
\label{sec:web-radial-matrix}
\hypertarget{proof:web-radial-matrix:s1}{%
\subsection{A genuinely global polynomial object}\label{proof:web-radial-matrix:s1}}

Let \(E=S^-\otimes L\), \(\det E=L^3\), and use the actual bounded auxiliary section \(v\). Define \[
 \mathcal B=v\wedge\nabla v\in\Omega^1(L^3).
\] No division, moving frame, or lower-amplitude cutoff is needed for this definition. If \(L_E=\nabla^*\nabla\), then the exact identities are \[
 \delta_{3A}\mathcal B=v\wedge L_Ev,
 \tag{\thesection.1}\label{proof:web-radial-matrix:e1}
\] \[
 (d_{3A}\mathcal B)_{ij}
 =2\nabla_i v\wedge\nabla_jv+v\wedge F^E_{ij}v.
 \tag{\thesection.2}\label{proof:web-radial-matrix:e2}
\] The repeated derivative term in \eqref{proof:web-radial-matrix:e1} vanishes because each \(\nabla_i v\wedge\nabla_i v=0\). In \eqref{proof:web-radial-matrix:e2}, the exterior derivative is antisymmetric in \(i,j\).

The actual endpoint estimates therefore give \[
 \|\mathcal B\|_{2,\rm av}\le C,
 \qquad \|\delta_{3A}\mathcal B\|_{2,\rm av}\to0.
 \tag{\thesection.3}\label{proof:web-radial-matrix:e3}
\] The self-dual part of \(v\wedge F^Ev\) is \(L^2\)-small: the traceless self-dual block of the negative spin connection is controlled by \(\operatorname{Ric}^{\circ}\), and the central determinant-curvature term wedges with \(v\) to zero. The remaining quadratic term in \eqref{proof:web-radial-matrix:e2} is only uniformly \(L^1\) under the unrestricted hypotheses. Equations \eqref{proof:web-radial-matrix:e1}–\eqref{proof:web-radial-matrix:e3} are actual global statements, not formal limiting equations.

\hypertarget{proof:web-radial-matrix:s2}{%
\subsection{Local exact model retaining the radial variables}\label{proof:web-radial-matrix:s2}}

The following local formulas are the exact zero-error first-jet calculation. They must not be used to discard derivative errors in an actual weak sequence without testing them.

Use \(\rho=\rho_*=4/3\), \(m=|v|^2>0\), and set \[
 u={dm\over2m},\quad z=u_1-iu_2,\quad
 q=(u_2,-u_1,u_4,-u_3),\quad d=-(u_3+iu_4).
\] In the adapted frame, the complete phase/Dirac relations give \[
 b=d\theta_1+c\theta_2,
 \qquad \theta_1=e^1-ie^2,\quad\theta_2=e^3-ie^4,
 \qquad \mathcal B=mb.
 \tag{\thesection.4}\label{proof:web-radial-matrix:e4}
\] Thus \(d\) is not set to zero. Write \[
 \Psi_1=e^{12}+e^{34},\quad
 \Xi_+=e^{13}-e^{24}+i(e^{14}+e^{23}),
\] so \(\overline{\Xi_+}=\theta_1\wedge\theta_2\), and both are self-dual. The exact Cartan formulas are \[
 d\theta_1=-i(\mathfrak a-q)\wedge\theta_1
                 +\kappa\Xi_+-d\overline{\Xi_+},
\] \[
 d\theta_2=i(2\mathfrak a-q)\wedge\theta_2
                 -\bar c\Xi_- -2i\bar d\,e^{12},
 \qquad \kappa={\sqrt m\over2\rho_*^{3/2}}.
\] Here \(A=i\mathfrak a\), and \(\Xi_-\) is the anti-self-dual form used in the preceding continuation.

Since the off-diagonal connection curvature is \(F^E_{21}=d_{3A}b-2iq\wedge b\), \eqref{proof:web-radial-matrix:e2} becomes \[
 d_{3A}^+\mathcal B
 =m(F^E_{21})^++2m\bigl((u+iq)\wedge b\bigr)^+.
\] The elementary wedge computation gives \[
 \boxed{d_{3A}^+\mathcal B
 =m(F^E_{21})^++2m(d^2+c\bar z)\overline{\Xi_+}.}
 \tag{\thesection.5}\label{proof:web-radial-matrix:e5}
\] Also \(\delta_{3A}b=2u\mathbin{\lrcorner}b\) is precisely the off-diagonal component of \(L_Ev=2v\), so \eqref{proof:web-radial-matrix:e1} yields \(\delta_{3A}\mathcal B=0\) in this zero-error model.

The complete off-diagonal self-dual curvature equation itself is \[
 \begin{aligned}
 (F^E_{21})^+={}&\mathcal L_1(D_2d-iqd)
              +\mathcal L_2(D_5c-3iqc)\\
 &+\kappa d\Xi_+-d^2\overline{\Xi_+}
              -ic\bar d\Psi_1,
 \end{aligned}
 \tag{\thesection.6}\label{proof:web-radial-matrix:e6}
\] where \(\mathcal L_a(\alpha)=(\alpha\wedge\theta_a)^+\). Thus the \(d\)-equation and \(c\)-equation are coupled; the original fifth-line CR equation results only after a justified removal of the radial terms.

For example, taking \(D=m^{-3/2}d\), \(C=m^{-3/2}c\) cancels the \(d^2\overline{\Xi_+}\) terms but leaves \[
 m^{-3/2}(F^E_{21})^+
 =\mathcal L_1(D_2D)+\mathcal L_2(D_5C)
       +\kappa D\Xi_+ +2i(C\bar d-Dz)\Psi_1.
 \tag{\thesection.7}\label{proof:web-radial-matrix:e7}
\] The last radial transport term does not vanish. The expressions with negative powers of \(m\) in \eqref{proof:web-radial-matrix:e7} are only local on \(m>0\); they are not an independently global replacement for \(\mathcal B\).

\hypertarget{proof:web-radial-matrix:s3}{%
\subsection{The integrated operator has one negative potential direction}\label{proof:web-radial-matrix:s3}}

For a unitary line bundle \(K\), write \(F=-iF^K\) for its real curvature and use its matrix action \((F\alpha)_j=F_{jk}\alpha_k\). On line-valued one-forms, the ASD deformation Weitzenböck formula is \[
 \delta_K^*\delta_K+2(d_K^+)^*d_K^+
 =\nabla_K^*\nabla_K+\operatorname{Ric}-2iF^-.
 \tag{\thesection.8}\label{proof:web-radial-matrix:e8}
\] It follows by adding \(d_K^**d_K\) to the Hodge Laplacian: the line-curvature terms are respectively \(+i*F\) and \(-iF\), hence \(-2iF^-\). Only Ricci, not the full tangent curvature, appears in \eqref{proof:web-radial-matrix:e8}.

In the endpoint model the real curvature of \(L^3\) has \[
 F_{L^3}^-=\Omega_-=e^{12}-e^{34},\qquad\operatorname{Ric}=-g.
\] Since \(\Omega_-\theta_1=-i\theta_1\) and \(\Omega_-\theta_2=i\theta_2\), the zero-order operator in \eqref{proof:web-radial-matrix:e8} has eigenvalues \[
 \boxed{-3\text{ on }\theta_1,\qquad +1\text{ on }\theta_2.}
 \tag{\thesection.9}\label{proof:web-radial-matrix:e9}
\] Consequently, the integrated identity has the explicit model form \[
 \int\bigl(|\delta_{3A}\mathcal B|^2+2|d_{3A}^+\mathcal B|^2\bigr)
 =\int\bigl(|\nabla\mathcal B|^2
           -6m^2|d|^2+2m^2|c|^2\bigr).
 \tag{\thesection.10}\label{proof:web-radial-matrix:e10}
\] For a genuine approximate sequence, retain the additional exact curvature-model error pairing \[
 \int\operatorname{Re}\langle
   (\operatorname{Ric}+g-2i(F_{L^3}^--\Omega_-))\mathcal B,
        \mathcal B\rangle.
 \tag{\thesection.11}\label{proof:web-radial-matrix:e11}
\] That error is not automatically small: the coefficient is \(L^2\)-small, but \(\mathcal B\) is presently only bounded in \(L^2\). An \(L^4\) bound or a separately justified bounded-test formulation would be required to discard \eqref{proof:web-radial-matrix:e11}.

In \eqref{proof:web-radial-matrix:e5}, \(|\overline{\Xi_+}|^2=4\). Hence with vanishing curvature error, \[
 2|d_{3A}^+\mathcal B|^2=32m^2|d^2+c\bar z|^2.
 \tag{\thesection.12}\label{proof:web-radial-matrix:e12}
\] The left side is thus a radial quartic term; setting it equal to zero before controlling the radial variables would assume the missing conclusion.

\hypertarget{proof:web-radial-matrix:s4}{%
\subsection{The unavoidable normal derivative energy is insufficient for positivity}\label{proof:web-radial-matrix:s4}}

Because \(\mathcal B\) is of the \(\theta_1,\theta_2\) type, its covariant derivative has an opposite-type component determined by the positive spinor connection. That component cannot be cancelled by differentiating \(d,c\).

The Clifford-to-Cartan calculation gives \[
 (\nabla\theta_1)^{\rm opp}=\kappa\bar\theta_1\otimes\bar\theta_2,
 \qquad
 (\nabla\theta_2)^{\rm opp}=-\kappa\bar\theta_1\otimes\bar\theta_1.
\] Since both \(\theta_a\) have squared norm two, \[
 \boxed{|(\nabla\mathcal B)^{\rm opp}|^2
       =4m^2\kappa^2(|d|^2+|c|^2).}
 \tag{\thesection.13}\label{proof:web-radial-matrix:e13}
\] For variable \(m\), adding this forced normal part to the potential gives exactly \[m^2(m/\rho_*^3-6)|d|^2+m^2(m/\rho_*^3+2)|c|^2.\] At the frozen zeroth-order value \(m=m_*\), \(\kappa^2=1/6\), so the coefficients become \[
 -{16\over3}m^2|d|^2+{8\over3}m^2|c|^2.
 \tag{\thesection.14}\label{proof:web-radial-matrix:e14}
\] In particular the radial coefficient remains negative throughout the lower-well range \(0<m\le m_*\). Freezing this coefficient is only a sign check, not an assertion that a variable-amplitude sequence has \(dm=0\), and no coefficient replacement under an uncontrolled derivative-energy integral is used. It therefore does not yield coercivity of the complete radial system. This coefficient calculation directly rules out the proposed shortcut in which the forced normal energy alone pays the entire negative radial potential.

\hypertarget{proof:web-radial-matrix:s5}{%
\subsection{Combining with the scalar mass equation does not yet supply the missing sign}\label{proof:web-radial-matrix:s5}}

Let \(E_1=m_1^2+m_2^2\), \(E_2=m_3^2+m_4^2\). The same complete first-jet system gives \[
 |\nabla v|^2={E_1+2E_2\over2m}+2m|c|^2.
\] Integrating \(\Delta m=4m-2|\nabla v|^2\) against \(m\), with the positive Laplacian convention, yields \[
 2\int E_1+3\int E_2+4\int m^2|c|^2=4\int m^2.
 \tag{\thesection.15}\label{proof:web-radial-matrix:e15}
\] Also \(6m^2|d|^2=(3/2)E_2\). Eliminating the \(c\)-mass from the potential in \eqref{proof:web-radial-matrix:e10} using \eqref{proof:web-radial-matrix:e15} gives \[
 \int(-6m^2|d|^2+2m^2|c|^2)
 =\int(2m^2-E_1-3E_2).
 \tag{\thesection.16}\label{proof:web-radial-matrix:e16}
\] The radial energies in \eqref{proof:web-radial-matrix:e16} have the wrong sign for an immediate positive lower bound. The quartic term \eqref{proof:web-radial-matrix:e12}, the remaining same-type derivatives of \(\mathcal B\), and the actual error pairing \eqref{proof:web-radial-matrix:e11} have not been controlled by a new compensating inequality. No sign has been assigned to them without proof.


\section{The exact positive-spinor Hessian limit}
\label{sec:web-exact-hessian}
The exact identity is \[
\int|H|^2=\int|L\Phi|^2
-\int\operatorname{Ric}_{ij}\operatorname{Re}\langle P_i,P_j\rangle
+\sum_{ij}\int\operatorname{Re}\langle F_{ij}P_j,P_i\rangle
+\tfrac12\sum_{ij}\int|F_{ij}\Phi|^2.
\tag{\thesection.1}\label{proof:web-exact-hessian:e1}
\]

At the calibrated endpoint choose the compatible frame \[
\Phi=\sqrt{4/3}\,e_1,\qquad
P_1=t e_2,\quad P_2=-it e_2,\quad P_3=P_4=0,
\qquad t^2=2/9.
\] The real determinant curvature is \[
-iF_A=(2/3)(e^{12}+e^{34})+(1/3)(e^{12}-e^{34})
=e^{12}+(1/3)e^{34}.
\] With the present setting's Clifford signs, the positive-spinor curvature is \[
F_{ij}=\tfrac16c^+(e^{ij})+\tfrac12(F_A)_{ij}I.
\] In particular \(F_{12}=\operatorname{diag}(i/3,2i/3)\) and \(F_{34}=\operatorname{diag}(0,i/3)\); each of the four mixed pairs contributes \(\rho/36\) to \(|F_{ij}\Phi|^2\).

The four terms in \eqref{proof:web-exact-hessian:e1} consequently have the following limits, in their displayed order: \[
\frac4{27},\qquad\frac{12}{27},\qquad
\frac8{27},\qquad\frac8{27}.
\] Thus \[
\boxed{\fint|\nabla_A^2\Phi|^2\longrightarrow\frac{32}{27}.}
\tag{\thesection.2}\label{proof:web-exact-hessian:e2}
\] The passage to the limit in \eqref{proof:web-exact-hessian:e1} uses the bounded-model-plus-small-\(L^2\) curvature fit and strong \(L^2\) convergence of \(|P|^2\); no differentiated curvature or \(W^-\) budget is used.

To see what this pays, write \(H_{ji}=(h_{ji},k_{ji})\). In the calibrated first jet, \(\nabla\zeta=0\) fixes \[
h_{ji}=-\overline{\alpha_j}\alpha_i/\sqrt\rho,
\qquad (\alpha_1,\alpha_2,\alpha_3,\alpha_4)=(t,-it,0,0),
\] so \(|h|^2=4/27\). Differentiated Dirac gives \[
k_{j1}-ik_{j2}=h_{j3}+ih_{j4},\qquad
k_{j3}-ik_{j4}=-h_{j1}-ih_{j2}.
\] The complementary squared sum on the right is \(8/27\). The pure-kernel quadratic contribution to \(\nabla\tau\) vanishes, leaving \[
\nabla_j\tau=\rho(k_{j1}+ik_{j2},-k_{j3}-ik_{j4}).
\] Hence the same total Hessian budget yields \[
\begin{aligned}
\fint|\nabla\tau|^2
&=\rho_*^2\left[2\left(\frac{32}{27}-\frac4{27}\right)-\frac8{27}\right]+o(1)\\
&=\frac{256}{81}+o(1)=2m_*+o(1).
\end{aligned}
\tag{\thesection.3}\label{proof:web-exact-hessian:e3}
\] This is exactly the total torsion-gradient budget already obtained from its Lichnerowicz formula. At the calibrated algebraic level it does not assign a separate positive excess to the radial part of \(\nabla\tau\).

Coefficient replacement in the limiting decomposition must not be justified merely by bounded Hessian energy. It can instead be localized with a weight \((\rho-\rho_*)^2\) in the positive-spinor Hessian interchange; the weight-derivative cost is bounded by \(C\int|d\rho|^2|P|^2=o(V)\), already available from the strong-current estimate. This removes an amplitude-coefficient concentration ambiguity; it does not create a new radial term in \eqref{proof:web-exact-hessian:e3}.

\textbf{Conclusion:} the exact positive-spinor Hessian identity does sharpen the known coarse budget to \(32/27\), but its calibrated torsion content is the existing \(2m_*\) total-gradient identity. It has not supplied a new inequality forcing \(d|\tau|^2\) or \(d|v|^2\) to vanish. Counting \eqref{proof:web-exact-hessian:e2} and \eqref{proof:web-exact-hessian:e3} as two independent coercive budgets would double-count the same information.


\section{Why the rank-one lemma retains a Ricci budget}
\label{sec:web-line-limitation}
\hypertarget{proof:web-line-limitation:Claim-that-cannot-be-used}{%
\subsection{Claim that cannot be used}\label{proof:web-line-limitation:Claim-that-cannot-be-used}}

The following properties, without any control of the geometry or an equation for the section, do \textbf{not} force normalized first-Chern square to vanish:

\[
 |\sigma_j|\to1\text{ in measure},\qquad
 \|\sigma_j\|_{L^4_{\rm avg}}+\|\nabla\sigma_j\|_{L^2_{\rm avg}}\le C,
 \qquad |F_j|\le C.
\tag{\thesection.1}\label{proof:web-line-limitation:e1}
\]

In particular, replacing the last condition by uniform integrability of curvature squares does not make the implication valid.

\hypertarget{proof:web-line-limitation:Construction}{%
\subsection{Construction}\label{proof:web-line-limitation:Construction}}

Let L be a degree-one Hermitian line bundle on an oriented two-sphere. Choose three nested coordinate discs

\[
 D_0\Subset D_1\Subset D_2\Subset S^2.
\]

Choose a smooth section s with one simple zero in D0 and norm identically one outside D0. Such a section exists: the degree-one clutching description gives a single-index-one section, and its positive norm away from the zero can be rescaled to one outside D0. The Hermitian metric and the section are fixed from now on.

Choose a unitary connection whose curvature is supported outside D2. To see this directly, start with any unitary connection and choose a smooth imaginary two-form F supported outside D2 with

\[
 \frac{i}{2\pi}\int_{S^2}F=1.
\]

The difference between F and the original curvature is exact, so adding a global imaginary one-form to the connection realizes F. On D2 the connection is flat; since D2 is simply connected it admits a parallel frame there. These choices do not require s to be parallel.

Let g0 be a fixed smooth metric, round on D1. Choose a smooth function \(0\le\chi\le1\), equal to one on D1 and zero outside D2. For \(0<\epsilon<1\), set

\[
 g_\epsilon=\epsilon^{2\chi}g_0.
\tag{\thesection.2}\label{proof:web-line-limitation:e2}
\]

The curvature F has uniformly bounded norm with respect to these metrics, because the metrics equal g0 on its support. Their areas converge to a positive number: the region outside D2 keeps its fixed positive area.

The amplitude defect is supported in D0, whose area is multiplied by \(\epsilon^2\). Thus \(|s|\to1\) in normalized measure. On the other hand, the covariant Dirichlet energy in two dimensions is conformally invariant:

\[
 \int_{S^2}|\nabla s|_{g_\epsilon}^2\,dA_{g_\epsilon}
 =\int_{S^2}|\nabla s|_{g_0}^2\,dA_{g_0}=:E_s<\infty.
\tag{\thesection.3}\label{proof:web-line-limitation:e3}
\]

Now take

\[
 M=S^2\times S^2,
 \qquad G_\epsilon=g_\epsilon\oplus g_\epsilon,
 \qquad E=\operatorname{pr}_1^*L\otimes\operatorname{pr}_2^*L,
 \qquad \sigma=s\otimes s.
\]

Use the product connection. The amplitude defect is supported in the union of the two product sets with one factor in D0, and hence has vanishing relative volume. The section is uniformly bounded. Its Dirichlet energy is

\[
 \int_M|\nabla\sigma|^2\,dV_{G_\epsilon}
 =2E_s\int_{S^2}|s|^2\,dA_{g_\epsilon},
\tag{\thesection.4}\label{proof:web-line-limitation:e4}
\]

which is bounded. The total volume is bounded below, so the normalized energy is bounded as well.

The curvature is the sum of the two pulled-back curvatures. Orthogonality of the product factors and the support property give a uniform pointwise curvature bound. Nevertheless,

\[
 \boxed{c_1(E)^2[M]=2,\qquad
 \frac{c_1(E)^2[M]}{\operatorname{Vol}(M,G_\epsilon)}
 \longrightarrow\frac2{V_0}>0.}
\tag{\thesection.5}\label{proof:web-line-limitation:e5}
\]

All properties \eqref{proof:web-line-limitation:e1} hold, and the claimed Chern-square conclusion fails.

\hypertarget{proof:web-line-limitation:Why-this-does-not-contradict-the weak endpoint problem}{%
\subsection{Why this does not contradict the weak endpoint problem}\label{proof:web-line-limitation:Why-this-does-not-contradict-the weak endpoint problem}}

The construction does not have nearly negative Einstein Ricci curvature. On D1, round curvature rescales by \(\epsilon^{-2}\), and its squared integral diverges. The covariant rough Laplacian of s also need not be uniformly controlled.

Therefore a rank-one argument for the actual derived section \(\sigma\in L^5\) must use additional information from the original system. The most natural candidates are the normalized small Einstein residual, a genuine equation for sigma, or a signed differential identity linking its phase current to the actual self-dual curvature. A proof using only \eqref{proof:web-line-limitation:e1}, even with the pointwise curvature strengthening shown here, cannot work.

\hypertarget{proof:web-line-limitation:The-same-counterexample-can-have-harmonic-curvature}{%
\subsection{The same counterexample can have harmonic curvature}\label{proof:web-line-limitation:The-same-counterexample-can-have-harmonic-curvature}}

The failure is not repaired merely by adding \(\delta F=0\). On each sphere replace the connection by one with harmonic curvature

\[
 F_\epsilon=-\frac{2\pi i}{\operatorname{Area}(g_\epsilon)}
             dA_{g_\epsilon}.
\]

These forms have uniformly bounded pointwise norms in the varying metrics. They also have uniformly bounded coefficients in a fixed g0 norm, because the conformal factors in \eqref{proof:web-line-limitation:e2} are at most one and the areas stay bounded below.

Fix one reference connection with curvature \(F_*\). On the fixed sphere, choose its connection difference as

\[
 \alpha_\epsilon=d_{g_0}^*\Delta_{g_0}^{-1}(F_\epsilon-F_*).
\]

The difference has integral zero, and therefore \(d\alpha_\epsilon=F_\epsilon-F_*\). Fixed-metric elliptic estimates give a uniform W1,p bound for every finite p, hence a uniform C0 bound for p greater than two. The fixed section s retains a uniformly bounded covariant Dirichlet energy by two-dimensional conformal invariance and this connection bound.

On the product of the two identical spheres, the sum of the two pulled-back harmonic curvatures is self-dual and parallel. The amplitude, energy, and Chern-square conclusions above still hold. Ricci curvature remains uncontrolled. This strengthened version isolates why the Ricci hypothesis in the logarithmic resolvent lemma is substantive rather than optional.


\section{Orthogonal obstruction estimates and their scale}
\label{sec:web-vector-scope}
\hypertarget{proof:web-vector-scope:s2}{%
\subsection{The precise vector lemma}\label{proof:web-vector-scope:s2}}

Let H be a real Hilbert space, e a unit vector, and write u = x e + y with y perpendicular to e. Suppose

\[
 |x|\le\epsilon\|y\|+r,
 \qquad \epsilon,r\ge0.
\tag{\thesection.1}\label{proof:web-vector-scope:e1}
\]

For any real a,

\[
 \boxed{\|a e+u\|\ge
 \frac{(|a|-r)_+}{\sqrt{1+\epsilon^2}}.}
\tag{\thesection.2}\label{proof:web-vector-scope:e2}
\]

Indeed,

\[
 |a|\le |a+x|+|x|
 \le |a+x|+\epsilon\|y\|+r
 \le \sqrt{1+\epsilon^2}\|a e+u\|+r.
\]

For r = 0 this is sharp: the minimizing correction has

\[
 x=-\frac{a\epsilon^2}{1+\epsilon^2},
 \qquad \|y\|=\frac{|a|\epsilon}{1+\epsilon^2}.
\]

If an additional approximation error z is present, the lower bound loses at most \(\|z\|\). This is an actual coercive estimate, but its application requires proof of \eqref{proof:web-vector-scope:e1}. Decomposing an obstruction into parallel and perpendicular components alone supplies no such relation.

There is also a dimension issue. A family of relations

\[
 |x|\le\sum_{k=1}^{N_j}\epsilon_{jk}|y_k|
\]

gives \eqref{proof:web-vector-scope:e1} with \(\epsilon_j=(\sum_k\epsilon_{jk}^2)^{1/2}\). Knowing that every coefficient is small is insufficient when \(N_j\to\infty\). For instance \(\epsilon_{jk}=N_j^{-1/2}\) tends to zero individually, but its operator norm remains one. An unrestricted the weak endpoint problem argument needs the operator-norm version, or an independent bound on the total obstruction dimension.

\hypertarget{proof:web-vector-scope:s3}{%
\subsection{The scale issue survives even after granting a bubble tree}\label{proof:web-vector-scope:s3}}

Suppose hypothetically that a weak the weak endpoint problem sequence has a suitable gluing description and that \eqref{proof:web-vector-scope:e2} produces a residual of size \(c t_j\), where \(t_j\downarrow0\) is a gluing parameter. Such a bound excludes an exact solution, and excludes an error \(o(t_j)\). It does not exclude an error that merely tends to zero.

The elementary scaling example is already decisive. A tensor of pointwise norm of order one supported in a region of volume of order \(t_j^2\) has squared L2 norm of order \(t_j^2\). Hence its normalized squared L2 norm tends to zero when the total volume stays bounded below. A nonzero obstruction on every shrinking core can therefore be consistent with a vanishing average residual.

This is not a proposed actual the weak endpoint problem counterexample: the monopole system and all matching conditions have not been imposed. It proves a logical point about the proposed import. A finite-scale nonvanishing theorem does not imply the volume-scale lower bound required by the weak endpoint problem. A successful adaptation must show that the aggregate obstruction has squared norm comparable to V, or that the actual field controls the residual relative to every relevant gluing scale. Neither follows from \(\fint|\mathrm{Ric}+g|^2\to0\).

\hypertarget{proof:web-vector-scope:s4}{%
\subsection{A global vector decomposition already available in the weak endpoint problem}\label{proof:web-vector-scope:s4}}

Assume an actual zero-endpoint sequence, scalar curvature minus four, and a fixed finite bound on

\[
 b_j=\fint|W_j^-|^2.
\]

Use the bounded resolvent section v from the preceding continuation and set

\[
 a=|v|^2,\quad m_*=128/81,\quad Q=-iq_-(v),\quad q_0^2=m_*^2/8.
\]

The already established strong-current and Hopf-form calculations give

\[
 a\to m_*\text{ in }L^2_{\rm avg},\qquad |v|\le C,
\]

\[
 \fint\langle W^-Q,Q\rangle\to5m_*^2/12,
 \qquad |Q|^2=a^2/8.
\tag{\thesection.3}\label{proof:web-vector-scope:e3}
\]

Define the smooth polynomial tensor

\[
 W_0=\frac5{q_0^2}
 \left(Q\otimes Q-\frac{|Q|^2}{3}I\right).
\tag{\thesection.4}\label{proof:web-vector-scope:e4}
\]

No normalization by Q is used, so \eqref{proof:web-vector-scope:e4} remains smooth at its zeros. Since W is trace free,

\[
 \fint\langle W^-,W_0\rangle
 =\frac5{q_0^2}\fint\langle W^-Q,Q\rangle\to50/3.
\]

Also

\[
 |Q\otimes Q-|Q|^2I/3|^2=\frac23|Q|^4,
 \qquad \fint|W_0|^2\to50/3.
\]

Writing R = W minus W0, one obtains, for \textbf{every} bounded-b zero sequence,

\[
 \boxed{\fint\langle R,W_0\rangle\to0,
 \qquad \fint|R|^2=b_j-50/3+o(1).}
\tag{\thesection.5}\label{proof:web-vector-scope:e5}
\]

The previous endpoint theorem used the special case where the right side tends to zero. Formula \eqref{proof:web-vector-scope:e5} records what remains when it does not. It is a restatement of that established pairing, not a new threshold theorem.

\hypertarget{proof:web-vector-scope:s4.1}{%
\subsubsection{Which vector components can carry the excess?}\label{proof:web-vector-scope:s4.1}}
